\section{Fields with controlled Selmer groups}\label{sec:fields}

The connectedness argument needs large totally real fields, very few allowed
non-dyadic primes, and a strict Selmer bound at a small proportion of the
dyadic places.  We construct these fields here.  The original allowed set,
including the primes needed for the modular seed, is fixed before this
construction.

Let $\Sigma$ be a finite set of odd rational primes, and let $\bar r$ be a
semisimple representation of $G_{\QQ}$ over a sufficiently large finite
field $k$ of characteristic two, unramified outside $\Sigma\cup\{2\}$.
If $\bar r$ is reducible, write its two characters as
$\bar\chi_1,\bar\chi_2$ and put $\alpha=\bar\chi_1/\bar\chi_2$.
Thus $\alpha$ has odd order, is even, and is unramified at $2$.
The last assertion follows from
$\ZZ_2^\times$ being a pro-$2$ group.  Write $\alpha_2$ for its
restriction to $G_{\QQ_2}$.
For a number field $E$ and a finite set $S_E$ containing the dyadic and
real places, all Selmer groups below are subspaces of
$H^1(G_{E,S_E},k(\beta))$; at places outside $S_E$ their local condition
is unramified.

\begin{proposition}[Field preparation]\label{prop:field-preparation}
There is a constant $U$, depending only on the fixed residual data and
$\Sigma$, with the following property.  One can construct abelian,
totally real $2$-extensions
\[
                 \QQ\subset A\subset F
\]
in which $2$ splits completely, together with allowed sets $S_A,S_F$,
such that the following assertions hold.
\begin{enumerate}
\item The allowed sets contain the real and dyadic places and all places
above $\Sigma$.  Their numbers $u_A,u_F$ of non-dyadic finite places
are at most $U$.  The degrees $d_A=[A:\QQ]$ and $[F:A]$ can be
made arbitrarily large, with $d=[F:\QQ]=[F:A]d_A$.
\item If $\bar r$ is absolutely irreducible, its restriction to $G_F$
is absolutely irreducible.  In this case there are no marked places.
\item In the reducible case, the characters remain distinct on $G_F$
if they were distinct on $G_{\QQ}$.  If $\alpha_2\ne1$, there are
again no marked places.  If $\alpha_2=1$, there is a dyadic place
$v_0$ of $A$ and the marked set $M$ consists of all its extensions
to $F$.  Set $m=|M|$; thus $m=[F:A]$ and $m/d=1/d_A$ in the
marked case, whereas $m=0$ otherwise.
\item In the reducible case, for $\beta=\alpha,\alpha^{-1}$ the
following group is zero:
\begin{equation}\label{eq:dual-vanishing}
 \left\{c\in H^1(G_{F,S_F},k(\beta)):
       \operatorname{loc}_v(c)=0\text{ for every }v\in S_F\setminus M
 \right\}=0.
\end{equation}
Here the zero conditions include all real places.  The condition at a
marked place is unrestricted.
\item In the reducible case, restriction preserves residual extension
classes: for either ratio
$\beta$, the map
\[
 H^1(G_{A,S_A},k(\beta))\longrightarrow H^1(G_{F,S_F},k(\beta))
\]
is injective.
\end{enumerate}
In the reducible case let $H^1_{\mathrm{str}}(F,k(\beta))$ denote the
Selmer group with zero condition at $M$, full condition at
$S_F\setminus M$ (including the real places), and unramified condition
elsewhere.  Then
\begin{equation}\label{eq:strict-selmer}
 \dim_k H^1_{\mathrm{str}}(F,k(\beta))\le
 \begin{cases}
 d+u_F,&\alpha_2\ne1,\\
 2d-3m+u_F,&\alpha_2=1.
 \end{cases}
\end{equation}
\end{proposition}

The field construction has three cases.  For absolutely irreducible
residue, a cyclic extension disjoint from the residual splitting field
will suffice.  For scalar residue, a narrow-class-group argument will
give the required vanishing.  For distinct residual characters, we will
choose conductor primes successively so that changing the local Selmer
conditions removes classes without introducing new ones.  In each case
we must bound the number of allowed primes independently of the degrees
of the fields; this is what permits the strict dimension inequalities
used in Section~\ref{sec:connectedness}.

We use class field theory with real places in the modulus
\cite[Chapter~V, Theorems~1.7, 3.5, and~3.6]{MilneCFT}, the Kummer
identification, and inflation--restriction.  Local and global duality
give the following Selmer dimension formula for a finite coefficient
module $V$ and linear local conditions $\mathcal L$:
\begin{equation}\label{eq:field-selmer-formula}
 \begin{aligned}
 \dim H^1_{\mathcal L}(V)-\dim H^1_{\mathcal L^\perp}(V^\vee(1))
 &=h^0(V)-h^0(V^\vee(1))\\
 &\quad+\sum_v\bigl(\dim\mathcal L_v-h^0_v(V)\bigr).
 \end{aligned}
\end{equation}
This is obtained by taking dimensions in Poitou--Tate and using the
global Euler formula
\cite[Chapter~I, Theorems~4.10 and~5.1, Remark~5.2(a)]{Milne2006}.
The real places are retained; their $h^0_v$ in this formula is ordinary
invariant dimension.  Our coefficient characters are trivial at real
places, so local $H^1$ there has dimension one.

\subsection{Cyclic fields and a scalar genus argument}

Fix integers $n>a\ge1$.  If an odd prime $r$ satisfies
\begin{equation}\label{eq:conductor-power}
 r\equiv1\pmod{2^{n+1}},\qquad
               2\in(\FF_r^\times)^{2^n},
\end{equation}
the degree-$2^n$ cyclic subfield $L(r,n)$ of $\QQ(\zeta_r)$ is real
and has $2$ split completely.  It is totally ramified at $r$ and
unramified at other finite primes.  For an odd prime $\ell\ne r$,
the Frobenius of $\ell$ generates its Galois group precisely when
$\ell$ is a nonsquare modulo $r$.  Indeed, both assertions are read
in the cyclic quotient of $\FF_r^\times$ of order $2^n$; an element
generates that quotient if and only if its exponent in
$\FF_r^\times$ is odd.

The power conditions in \eqref{eq:conductor-power} say that $r$ splits
in the Galois $2$-extension
\[
 E_n=\QQ(\zeta_{2^{n+1}},2^{1/2^n}).
\]
This extension is unramified at odd primes.  Consequently, independent
square-root conditions at specified odd primes are compatible with
splitting in $E_n$: any nontrivial product of their squareclasses is
ramified at an odd prime, whereas every quadratic subfield of $E_n$
is unramified there.  Chebotarev
\cite[Chapter~V, Theorem~3.23]{MilneCFT} therefore supplies $r$ satisfying
\eqref{eq:conductor-power} with every $\ell\in\Sigma$ nonsquare
modulo $r$.  One can avoid any further fixed finite set of primes.

This already proves the field assertion for absolutely irreducible
$\bar r$: take $F=L(r,n)$ and its degree-$2^a$ subfield for $A$.
Choose $r$ unramified in the residual splitting field.  Every nontrivial
subextension of $F/\QQ$ is ramified at $r$, so $F$ is disjoint from
that splitting field.  The residual image is therefore unchanged.
Each $\ell\in\Sigma$ has one prime above it at every cyclic layer.
The prime $r$ need not be allowed, since restricting a representation
does not introduce ramification at a prime where it was unramified.

This completes the construction for absolutely irreducible residue:
the residual image is unchanged and the number of allowed primes is
bounded independently of the cyclic degree.  For scalar residue we
also need the Selmer vanishing in \eqref{eq:dual-vanishing}.  We obtain
it from a different fixed base field and the following elementary
form of the one-prime genus argument.

\begin{lemma}\label{lem:one-prime-genus}
Suppose that $K/E$ is a quadratic extension of totally real fields
and that exactly one finite prime of $E$ ramifies.  If the narrow
class number of $E$ is odd, the narrow class number of $K$ is odd.
For every totally real field with odd narrow class number, each
totally positive unit is a square.
\end{lemma}

\begin{proof}
If the narrow class number of $K$ were even, the nonzero
$\FF_2$-space $\operatorname{Cl}^+(K)/2$ would have a nonzero
invariant linear functional for $\Gal(K/E)$.  Here we use the
elementary fact that a nonzero module of a finite $2$-group in
characteristic two has nonzero invariants, applied to the dual module.
Class field theory would give a quadratic extension $H/K$,
unramified at every finite prime, which is Galois over $E$.
The group $G=\Gal(H/E)$ has order four.  Above the unique finite prime
ramified in $K/E$, its inertia group $I$ has order two and maps
isomorphically to $\Gal(K/E)$.  It thus has trivial intersection with
$\Gal(H/K)$.  This excludes $G\simeq\ZZ/4\ZZ$, whose only subgroup
of order two is $\Gal(H/K)$.  Hence $G$ is the Klein four group.
All inertia groups above that prime equal $I$, since $G$ is abelian,
and all other finite inertia groups are trivial.  The quadratic
extension $H^I/E$ is consequently unramified at all finite primes,
contrary to oddness of the narrow class number of $E$.  Ramification
at infinity is allowed by the narrow class field and causes no
exception to this argument.

For the last assertion, the signature cokernel of the units is a
$2$-group embedded in the narrow class group, so it is zero.
For a totally real field of degree $e$, Dirichlet's unit theorem
\cite[Theorem~5.1]{MilneANT} gives
$|\mathcal O_E^\times/(\mathcal O_E^\times)^2|=2^e$.
The signature map from this group onto $\{\pm1\}^e$ is therefore
an isomorphism.  Its kernel is the group of totally positive units
modulo squares.
\end{proof}

Now suppose $\alpha=1$.  Put
\[
 B=\QQ(\sqrt{17}),\qquad
 \pi=\frac{5+\sqrt{17}}2.
\]
The Minkowski bound \cite[Theorem~4.3]{MilneANT}, here
$\sqrt{17}/2<3$, says that every ideal class of
$B$ has a representative of norm one or two.  Both primes above two
are principal, generated by $\pi$ and its conjugate, so the ordinary
class number is one.  The units $-1$ and $4+\sqrt{17}$ realize every
signature; the latter has norm $-1$.  Thus the narrow class number
is also one.

Choose $r$ as above, additionally inert in $B$, and put
\[
       F=B L(r,n),\qquad A=B L(r,n)_a,
\]
where the subscript denotes the degree-$2^a$ layer.  Inertness in
$B$ is the extra condition that $17$ be a nonsquare modulo $r$.
It is compatible with all the preceding conditions, including when
$17\in\Sigma$, in which case it repeats an existing condition.
Each quadratic step of $B L(r,n)/B$ has exactly one finite ramified
prime, the prime over $r$.  Lemma~\ref{lem:one-prime-genus} implies
that $F$ has odd narrow class number and that its totally positive
units are squares.  Take $S_A,S_F$ to contain only the prescribed
places above $\Sigma$, two, and infinity; in particular $r$ is not
allowed.  There are at most two primes above any $\ell\in\Sigma$.

Let $w_1$ be the dyadic prime of $B$ generated by $\pi$, and $w_2$
the other dyadic prime.  Both embeddings of $\pi$ into $\mathbb R$ are
positive.  At $w_2$ it is a unit congruent to seven modulo eight:
the odd root of $X^2-5X+2$ is seven modulo eight.  Thus $\pi$ is not
a square in $B_{w_2}=\QQ_2$.

Choose any $v_0$ of $A$ above $w_1$.  We prove a stronger version of
\eqref{eq:dual-vanishing}, allowing arbitrary localization at every
place of $F$ above $w_1$.  It suffices to work with $\FF_2$
coefficients.  By Kummer theory a class in this enlarged group is
represented by a totally positive $x\in F^\times$ modulo squares.
Its valuations are even outside the set $T$ of dyadic places above
$w_1$, and it is locally a square at all places above $w_2$.
The valuation-parity map
\begin{equation}\label{eq:scalar-parity}
             [x]\longmapsto (v(x)\bmod2)_{v\in T}
                 \quad\hbox{into }\FF_2^T
\end{equation}
is injective even on the larger space of totally positive squareclasses
with even valuations outside $T$.  Indeed, zero parity gives
$(x)=\mathfrak a^2$; odd ordinary class number makes $\mathfrak a$
principal, and division by a square leaves a totally positive unit,
which is a square by Lemma~\ref{lem:one-prime-genus}.

Our Selmer subspace is stable under the $2$-group $\Gal(F/B)$.
If nonzero, it has a nonzero invariant vector.  The action on $T$ is
transitive, so injectivity of \eqref{eq:scalar-parity} forces the
parity vector of that invariant to be the all-ones vector.
The unique totally positive squareclass with this parity is $[\pi]$:
it has valuation one at every place above $w_1$, and the injectivity
just proved gives uniqueness.  But $\pi$ is nonsquare at every place
above $w_2$, since these completions are still $\QQ_2$.
This contradiction proves the required vanishing, also after extending
coefficients from $\FF_2$ to $k$.

Finally, the kernel of restriction of scalar $H^1$ from $A$ to $F$
consists of characters factoring through the cyclic group $\Gal(F/A)$.
It is generated over $k$ by the unique nonzero $\FF_2$-valued
quadratic character, and every nonzero multiple is ramified above $r$.
As $r\notin S_A$, these characters are excluded.  This proves the
restriction assertion in the scalar case.

\subsection{Switching local conditions}

The scalar case is now complete.  With distinct residual characters,
we will replace the class-group argument by successive changes of local
conditions.  The next lemma explains why a switch removes classes when
the opposite Selmer group supplies enough local evaluations.

For the rest of the construction assume $\alpha\ne1$.  Put
$k_0=\FF_2(\alpha)$ and $q=|k_0|\ge4$.  Both $\alpha$ and
$\alpha^{-1}$ are defined over $k_0$; the desired conclusions over
$k$ follow by extension of scalars, since the continuous cochain
complexes and all the specified local conditions commute with this
finite scalar extension.  Write $K_\alpha$ for the cyclic,
totally real, odd-degree field cut out by $\alpha$.

This construction uses the viewpoint of Ramakrishna's auxiliary-prime
switching strategy~\cite{Ramakrishna2002}; see also the exposition in
Tang~\cite[Section~3A2, proof of Theorem~3.4]{Tang2019}.
The rank-one characteristic-two variant needed here is proved below,
rather than obtained by applying either reference's lifting theorem.

We first record the exact local switch used below.  At an odd place
where the coefficient character is trivial and the residue cardinality
is one modulo four, local $H^1$ is two-dimensional.  Its unramified
line $L_{\mathrm{ur}}$ is self-orthogonal.  Any line generated by a
ramified quadratic character is transverse to it and is also
self-orthogonal, for the pairing with the opposite coefficient.
Indeed, the mod-two Hilbert pairing is alternating when $-1$ is a
local square, and a line in a two-dimensional nondegenerate alternating
space is its own annihilator.

\begin{lemma}[Local switching]\label{lem:selmer-switch}
Let $X$ be a Selmer group over a number field $E$ with one-dimensional
coefficient $k_0(\beta)$, and let $Y$ be the Selmer group with coefficient
$k_0(\beta^{-1})$ and orthogonal local conditions.  Let $T$ be a finite
set of odd places of $E$ where the coefficient character is trivial
and the residue cardinality is one modulo four.  Suppose that $X$ has
the unramified local condition at every place in $T$, and that evaluation
of $Y$ on unramified Frobenius gives a surjection
\[
                 Y\longrightarrow k_0^T.
\]
Replacing the unramified lines at $T$ by arbitrary transverse ramified
lines changes $X$ to
\[
       \ker\bigl(X\longrightarrow k_0^T\bigr).
\]
In particular the switch introduces no new classes, and it strictly
decreases $X$ if its own evaluation is nonzero.
\end{lemma}

\begin{proof}
Let $X^{\mathrm{rel}}$ have full condition at $T$ and the old
conditions elsewhere.  The relevant part of the Poitou--Tate comparison
sequence is
\[
 0\longrightarrow X\longrightarrow X^{\mathrm{rel}}
 \longrightarrow\bigoplus_{v\in T} H^1_v/L_{\mathrm{ur},v}
 \longrightarrow Y^\vee.
\]
The last arrow is the transpose of localization on $Y$ under the
perfect pairings
$H^1_v/L_{\mathrm{ur},v}\times L_{\mathrm{ur},v}\to k_0$.
It is injective by the hypothesis, so $X^{\mathrm{rel}}=X$.
The new Selmer group is a subspace of this unchanged relaxed group.
Its localizations must belong both to the old unramified line and the
new transverse line, hence must be zero.  This gives the asserted
kernel exactly.
\end{proof}

\subsection{Frobenius evaluations with splitting conditions}

Take $B=\QQ$ if $\alpha_2\ne1$.  If $\alpha_2=1$, choose a prime
$q_0\notin\Sigma\cup\{2\}$ with
\[
           q_0\equiv1\pmod8,\qquad
           \alpha(\operatorname{Frob}_{q_0})\ne1,
\]
and take $B=\QQ(\sqrt{q_0})$.  Such a prime exists by Chebotarev,
since the odd-degree extension $K_\alpha$ is disjoint from
$\QQ(\zeta_8)$.  In this second case let $w_1,w_2$ be the two
dyadic places of $B$ and let $\sigma$ be its nontrivial automorphism.

All later conductor primes will satisfy \eqref{eq:conductor-power},
split in $K_\alpha$, and avoid the original primes and $q_0$.
The first conductor $r_1$ also splits in $B$, with every
$\ell\in\Sigma$ nonsquare modulo $r_1$.  Compatibility follows
from the independence of the odd squareclasses just used and from
the coprimality of $[K_\alpha:\QQ]$ with every auxiliary $2$-degree.
For each subsequent conductor $r$, require every
$\ell\in\Sigma$ and every earlier conductor to be a square modulo
$r$.  These conditions say that $r$ splits in a finite Galois
$2$-extension $E_2$ obtained by adjoining $E_n$ and those square roots.
In the quadratic-$B$ case, $E_2$ is unramified at $q_0$.
Consequently $B\cap E_2=\QQ$; the same is true of
$B\cap K_\alpha E_2$.  We can therefore impose either splitting or
inertness in $B$ at a new conductor.  In particular, $B$ is not included
among the fields required to split at the final inert step.

\begin{lemma}[Evaluation freedom]\label{lem:field-evaluations}
Fix finitely many cohomology classes over $B$ with coefficients
$k_0(\beta)$, where $\beta=\alpha$ or $\alpha^{-1}$, and impose
the preceding splitting conditions in $K_\alpha E_2$.
Evaluations use chosen places above the rational prime and chosen
Frobenius representatives in a common finite quotient through which
the coefficient actions and the finitely many cocycles factor.
Numerical prescriptions refer to these choices.
\begin{enumerate}
\item At primes split also in $B$, a Frobenius evaluation on any
$k_0$-linearly independent collection of classes can be prescribed
arbitrarily.  If $B$ is quadratic, the pair of evaluations on a
space $W$ at the two places above such a prime is given by one
functional restricted to $W$ and $\sigma W$.  It is uniformly
distributed as two independent functionals on $W$ whenever
$W\cap\sigma W=0$.
\item In the quadratic case, for primes inert in $B$ the evaluations
on $W$ are uniformly distributed in $W^\vee$ whenever
$1+\sigma$ is injective on $W$.  The evaluations on $c$ and
$\sigma c$ at the single place above such a prime are equal.
\item Collections for different Frobenius orbits of coefficient
characters can be prescribed independently.  If
$\alpha^{-1}=\alpha^{2^j}$, coefficient Frobenius identifies the
two collections semilinearly, and preserves all zero and rank
conditions used here.
\end{enumerate}
Every compatible prescription is realized, for suitable choices of
these places and representatives, at infinitely many rational primes.
Uniform distribution refers to the finite evaluation quotients used
for this choice; no analytic density estimate is required.
Chebotarev determines simultaneous conjugacy classes.  The applications
below use only rank and nonvanishing conditions, unchanged by the
resulting coefficient scalings and interchange of the two split places.
\end{lemma}

\begin{proof}
Set $N=G_{B K_\alpha E_2}$.  The finite quotient
\[
 \Gal(B K_\alpha E_2/B)=C\times P
\]
has $C=\Gal(K_\alpha/\QQ)$ of odd order and $P$ a $2$-group.
On $k_0(\beta)$ the group $C$ acts nontrivially.  Averaging over $C$
and using $k_0(\beta)^C=0$ gives
$H^i(C\times P,k_0(\beta))=0$ for all $i\ge0$.
Inflation--restriction thus makes restriction to $N$ injective.
On $N$ the coefficients are trivial, so cocycles are additive
homomorphisms and coboundaries are zero.

For independent classes $c_1,\ldots,c_t$, their joint evaluation map
$N\to k_0^t$ has additive image stable under multiplication by
$\beta(C)$.  These scalars generate $k_0$ over $\FF_2$, so its
image is a $k_0$-subspace.  If it were proper, a nonzero linear
functional annihilating it would give a nontrivial linear combination
of the $c_i$ restricting to zero on $N$, a contradiction.  Its image
is therefore all of $k_0^t$.  The paired assertion follows by applying
this to $W+\sigma W$; trivial intersection makes its two summands
independent.  The conjugation formula for cocycles identifies these
with the evaluations at the two conjugate places.

For the inert assertion choose $g\in G_{\QQ}$ acting nontrivially
on $B$ and trivially on $K_\alpha E_2$.  Then $g^2\in N$, and
$g$ represents $\sigma$ on cohomology.  Since $\beta(g)=1$, for
$h\in N$ the cocycle identity gives
\begin{equation}\label{eq:inert-evaluation}
 c((gh)^2)=c(g^2)+\operatorname{ev}_h((1+\sigma)c).
\end{equation}
The first term is a fixed affine offset.  The proven evaluation
surjectivity on $(1+\sigma)W$ shows that this expression is uniform
in $W^\vee$ if $1+\sigma$ is injective on $W$.  Also
$c(g^2)=(\sigma c)(g^2)$ and
$(1+\sigma)\sigma c=(1+\sigma)c$ in cohomology.  Evaluations on
$N$ do not see coboundaries, so \eqref{eq:inert-evaluation} gives
equality for $c$ and $\sigma c$, including the affine offset.

For different Frobenius orbits, regard these evaluation images as
$\FF_2[C]$-modules.  They have different simple isotypes; semisimplicity
of $\FF_2[C]$ makes their joint image the product of the individual
images.  In the remaining case the automorphism $x\mapsto x^{2^j}$
of $k_0$ carries cocycles with coefficient $\alpha$ to those with
coefficient $\alpha^{-1}$.  It commutes with restriction,
conjugation, and evaluation.  Every ramified line below is generated
by an $\FF_2$-valued quadratic character, so these lines, as well as
the zero, full, and unramified conditions, are preserved.  All
evaluation matrices are entrywise Frobenius transforms and have the
same ranks.

The finite collection of cocycles factors through a finite quotient
after adjoining their coefficient fields.  Chebotarev realizes the
conjugacy class of the chosen element of this quotient.  Choosing a
place in the corresponding finite Galois extension above the rational
prime then chooses a Frobenius representative;
one may choose the prescribed representative in that class.
Conjugacy scales coefficient evaluations and may interchange the
two split places, which preserves every rank and nonvanishing
condition in use.  Removing finitely many primes leaves infinitely
many choices.
\end{proof}

\subsection{The distinct-character construction}

Given the chosen conductors $r_1,\ldots,r_t$, take a diagonal cyclic
quotient of the product of the degree-$2^n$ fields $L(r_j,n)$:
if $\gamma_j:G_{\QQ}\to\ZZ/2^n\ZZ$ defines its $j$th factor,
use $\gamma=\sum_j c_j\gamma_j$ with every $c_j$ odd.  Let $L$
be the field of $\gamma$, let $L_a$ be its degree-$2^a$ layer, and
write
\[
           F=BL,\qquad A=BL_a,\qquad \delta=\gamma\bmod2.
\]
When $B$ is quadratic, choose any $v_0$ of $A$ above $w_1$ and
mark all its extensions to $F$.
The character $\delta$ is the quadratic sum of the conductor
characters.  Every $r_j$ is totally ramified in $L$, and $2$ is
split.  At each original $\ell\in\Sigma$, the first summand of
$\gamma(\operatorname{Frob}_\ell)$ is odd and all later summands
are even, so its Frobenius has order $2^n$.
At a previous conductor, adjoining a new conductor does not change
the local quadratic line $k_0\delta$: the new quadratic character
is unramified and trivial there.  This follows from the imposed
square conditions and quadratic reciprocity, since all conductors
are one modulo four.

The allowed sets now include primes above all the $r_j$, in addition
to the original allowed places.  They do not need to include $q_0$.
For $\beta=\alpha,\alpha^{-1}$ define $X_B^\beta$ as follows:
\begin{itemize}
\item its real conditions are zero;
\item if $B=\QQ$, its dyadic condition is zero; if $B$ is quadratic,
its dyadic conditions are full at $w_1$ and zero at $w_2$;
\item its other conditions are unramified, except above the $r_j$,
where they are the lines $k_0\delta$.
\end{itemize}
At a conductor the coefficient is trivial, since $r_j$ splits in
$K_\alpha$.  These definitions therefore make sense, including at
a possible inert conductor.

The remaining objective is to kill both spaces $X_B^\beta$ while
keeping the number of conductors bounded.  First we explain why this
vanishing gives the desired Selmer condition over the large field $F$;
we then construct the switches that achieve it.

Killing both spaces $X_B^\beta$ suffices for
\eqref{eq:dual-vanishing}.  To see this, in the quadratic case first
enlarge the possible marked set to all places above $w_1$.
The resulting Selmer space over $F$ is stable under the cyclic
$2$-group $\Gal(F/B)$.  If it were nonzero it would have a nonzero
invariant class.  Since $k_0(\beta)^{G_F}=0$,
inflation--restriction identifies these invariant classes with
classes over $B$.  At two and infinity the local conditions descend
unchanged, because these places split completely.  Away from the
conductors, $F/B$ is unramified, so a class whose restriction is
locally unramified has unramified localization downstairs.  This
applies both outside the allowed set and where localization was
required to vanish.  At a conductor
the local extension is cyclic and totally ramified of degree $2^n$,
and the kernel of local restriction is precisely the line of its
quadratic character, namely $k_0\delta$.  Thus the descended class
belongs to $X_B^\beta$, giving the asserted sufficiency.  For
$B=\QQ$ the same argument has no enlarged marked set.

\smallskip
\noindent\emph{The case $B=\QQ$.}
Suppose first that $B=\QQ$, so $\alpha_2\ne1$.  For $X_B^\beta$
the Selmer dimension formula gives
\begin{equation}\label{eq:q-selmer-difference}
      \dim X_B^\beta-\dim (X_B^\beta)^{\mathrm{dual}}=-1.
\end{equation}
Here the dual group has coefficient $\beta^{-1}$ and orthogonal
conditions.  The real zero condition contributes $-1$; the dyadic
zero condition contributes zero, since its local invariants vanish;
all other local conditions contribute zero.  Both global invariant
terms vanish.  Each opposite dual is therefore nonzero.

Append a conductor split in $K_\alpha E_2$ at which every opposite
dual has nonzero evaluation, and every nonzero $X_B^\beta$ also has
nonzero evaluation.  Such a choice exists by
Lemma~\ref{lem:field-evaluations}: per coefficient isotype there
are at most two required nonvanishings, and the failure probability
is at most $2/q<1$.  Different isotypes can be prescribed separately;
Frobenius-conjugate inverse characters impose identical conditions.
Lemma~\ref{lem:selmer-switch} says that neither space gains classes
and that each nonzero space decreases.  Repeating kills both spaces.

\smallskip
\noindent\emph{Preparing the quadratic case.}
It remains to carry out this argument when $B$ is quadratic, so
$\alpha_2=1$.  We first impose an auxiliary vanishing over $\QQ$.
Let $X_0^\beta$ have zero dyadic condition, full real condition, and
the same non-dyadic conditions as $X_B^\beta$ (unramified away from
the conductors, and $k_0\delta$ at them).  Its dimension minus that
of its opposite dual is again $-1$: the dyadic contribution is
$-1$, the real contribution is zero, and all other contributions
are zero.  Restriction
$H^1(G_{\QQ},k_0(\beta))\to H^1(G_B,k_0(\beta))$ is injective,
since $H^0(G_B,k_0(\beta))=0$.  Rational classes and their opposite
duals can therefore be included in the evaluation spaces over $B$
in Lemma~\ref{lem:field-evaluations}.  Use the preceding
one-coordinate switching argument to
kill both $X_0^\beta$, requiring all the added conductors to split
also in $B$.  Lemma~\ref{lem:field-evaluations} permits this extra
splitting.  Keep all the resulting local switches.

Let $Y_B^\beta$ be the group with coefficient $\beta$ whose
conditions are dual to those of $X_B^{\beta^{-1}}$.  Thus its real
conditions are full, and its dyadic conditions are zero at $w_1$
and full at $w_2$.  Its other conditions are the same as those of
$X_B^\beta$.  The dimension formula and conjugation give
\begin{equation}\label{eq:b-selmer-comparison}
 \dim Y_B^\beta=\dim X_B^{\beta^{-1}}+1,
       \qquad \sigma X_B^\beta\subset Y_B^\beta.
\end{equation}
Indeed, for $X_B$ the two dyadic contributions are $2$ and $-1$,
the two real contributions are $-1$ each, and the remaining
contributions vanish.  The inclusion follows from the local
conditions: conjugation exchanges the dyadic conditions and the
real zero conditions are contained in the full ones.

The following table summarizes the four groups used in the quadratic
case. Every row has coefficient $k_0(\beta)$, and ``full'' means the
full local $H^1$. The two entries for $B$ are ordered as $(w_1,w_2)$;
the single entry for $\QQ$ is the condition at $2$.
\begin{center}
\small
\begin{tabular}{@{}llcc@{}}
\toprule
Group & Base field & Real places & Dyadic places \\
\midrule
$X_B^\beta$ & $B$ & $0$ & $(\mathrm{full},0)$ \\
$Y_B^\beta$ & $B$ & $\mathrm{full}$ & $(0,\mathrm{full})$ \\
$X_0^\beta$ & $\QQ$ & $\mathrm{full}$ & $0$ \\
$(X_0^{\beta^{-1}})^{\mathrm{dual}}$ & $\QQ$ & $0$ & $\mathrm{full}$ \\
\bottomrule
\end{tabular}
\end{center}
At every selected conductor all four groups use the line $k_0\delta$,
and at every other non-dyadic finite place they use the unramified
condition. The rational
opposite dual in the last row will keep $X_0=0$ while the paired
conditions over $B$ are switched.

\begin{lemma}\label{lem:disjoint-conjugate-selmer}
As long as all conductors used so far split in $B$ and both
$X_0^\beta$ are zero, one has
\[
 W\cap\sigma W=0
        \quad\text{for }W=X_B^\beta\text{ and }W=Y_B^\beta.
\]
\end{lemma}

\begin{proof}
The intersection is stable under $\sigma$.  If nonzero it has a
nonzero invariant class, which descends to $\QQ$ because
$k_0(\beta)^{G_B}=0$.  Its dyadic localization is zero: intersecting
either pair of conjugate dyadic conditions imposes zero at both
$w_1$ and $w_2$.  At infinity it satisfies at least the full condition
defining $X_0^\beta$; this observation also covers $W=Y_B^\beta$.
At every conductor, splitting in $B$ makes descent of the rational
line $k_0\delta$ immediate.  At other primes unramified in $B/\QQ$,
the unramified condition descends by restriction to inertia.
There is no condition to check at $q_0$: the coefficient is
unramified with nontrivial Frobenius there, so local $H^0$ and
$H^2$ vanish and the local Euler formula gives $H^1=0$.
Thus the nonzero descended class lies in $X_0^\beta$, a contradiction.
\end{proof}

We have arranged $X_0^\beta=0$ for both coefficient characters.
The next split switches must preserve these vanishings while decreasing
the nonzero spaces $X_B^\beta$.  The preceding lemma, together with
Lemma~\ref{lem:field-evaluations}, supplies the freedom to choose
paired evaluations needed for each new split conductor.

\smallskip
\noindent\emph{Split switches.}
Now suppose both $X_B$ spaces are nonzero.  Include the restrictions
of the rational opposite-dual classes in the evaluation spaces just
described.  Choose a new split conductor
with the following three properties, for each coefficient isotype:
\begin{enumerate}
\item the paired evaluation on $Y_B^\beta$ has rank two;
\item the paired evaluation on $X_B^\beta$ is nonzero;
\item the evaluation on the opposite dual of the zero space
$X_0^{\beta^{-1}}$ is nonzero.
\end{enumerate}
The third dual has dimension one, by its dimension difference.
The first space has dimension at least two by
\eqref{eq:b-selmer-comparison}.  Lemmas~\ref{lem:field-evaluations}
and~\ref{lem:disjoint-conjugate-selmer} make the paired evaluations
on each of the first two spaces independent uniform functionals.
For $t=\dim Y_B^\beta\ge2$, the rank failure probability is
\[
  1-(1-q^{-t})(1-q^{1-t})\le q^{-1}+q^{-2}.
\]
The second failure probability is at most $q^{-2}$, and the third
is $q^{-1}$.  Their sum satisfies
\begin{equation}\label{eq:split-switch-probability}
            (q^{-1}+q^{-2})+q^{-2}+q^{-1}<1.
\end{equation}
No independence between these three failure events is needed for
this union bound.  Different character isotypes are independent;
conjugate inverse characters give identical requirements under
coefficient Frobenius.  Hence an admissible conductor exists.

Switching at its two places makes both $X_B$ spaces smaller and
introduces no new classes.  The third requirement has a separate
purpose: Lemma~\ref{lem:selmer-switch}, now over $\QQ$, preserves
$X_0^\beta=0$.  Thus Lemma~\ref{lem:disjoint-conjugate-selmer}
remains applicable at the next split step.  Continue until both
$X_B$ spaces vanish or exactly one remains.

\smallskip
\noindent\emph{The final inert switch.}
In the latter situation write $X_B^\beta\ne0$ and
$X_B^{\beta^{-1}}=0$.  Formula~\eqref{eq:b-selmer-comparison}
implies
\begin{equation}\label{eq:one-survivor}
 \dim X_B^\beta=\dim Y_B^\beta=1,
            \qquad Y_B^\beta=\sigma X_B^\beta.
\end{equation}
If the inverse characters were Frobenius conjugate, their $X_B$
spaces would have equal dimension, so this case could not arise.

Choose one final conductor inert in $B$, still split in
$K_\alpha E_2$.  Immediately before this step
Lemma~\ref{lem:disjoint-conjugate-selmer} holds, so $1+\sigma$ is
injective on each relevant $X_B$ and $Y_B$: its kernel would be a
fixed vector in $W\cap\sigma W$.
At this conductor require nonzero evaluation on both opposite duals
$Y_B^\alpha,Y_B^{\alpha^{-1}}$ and on the remaining $X_B^\beta$.
There is no third independent demand in the surviving isotype.
By \eqref{eq:one-survivor} and the exact equality in
Lemma~\ref{lem:field-evaluations}, the evaluations on $Y_B^\beta$
and $X_B^\beta$ coincide.  The affine offset in
\eqref{eq:inert-evaluation} does not change this equality.
Uniform evaluation and the union bound give the conservative failure
bound $2/q<1$ per isotype.  The switch at the single place over this
inert prime therefore kills the last $X_B$ space, while the opposite
zero space stays zero.  This is the final step: preservation of
$X_0=0$ after it is neither required nor asserted.  Its pre-step
vanishing was used precisely to justify injectivity of $1+\sigma$.

\subsection{Uniformity, restriction, and the numerical bounds}

The constructions now give the required dual Selmer vanishing.
To finish Proposition~\ref{prop:field-preparation}, we bound the number
of conductors, check preservation of residual data, and convert that
vanishing into the strict Selmer bounds.  The conductor count is the
point that allows both cyclic exponents to grow.

The number of conductors in the preceding argument is bounded
independently of $n$ and $a$.  Here is an explicit accounting which
also avoids making any uniformity assumption about their sizes.
Let $b=[B:\QQ]\le2$, and let
\[
 C_E=\sum_{\beta\in\{\alpha,\alpha^{-1}\}}
    \dim_{k_0}H^1(G_{E,T_E},k_0(\beta)),\qquad E=\QQ,B,
\]
where $T_E$ is the fixed set above $\Sigma$, two, and infinity.
In the quadratic case one may include $q_0$ in $T_B$ without change
to the relevant local cohomology.  Count both characters in this
sum even if they are Frobenius conjugate; the resulting bound is
only larger.  Allowing one new rational prime freely adds at most
$2b$ to each global $H^1$ dimension over $B$, by restriction to
inertia or by the localization sequence.  Thus it adds at most
$4b$ to their sum.  This deliberately generous bound applies
uniformly to every conductor and every quadratic line chosen there.

After the first conductor, the sum of the two dimensions over $\QQ$
is at most $C_{\QQ}+4$.  The first loop killing the auxiliary $X_0$
spaces therefore has length
\[
                         N_0\le C_{\QQ}+4.
\]
During it, the sum of the $B$-dimensions is at most
$C_B+4b(1+N_0)$.  Every later split switch decreases that sum by
at least one and never increases either space.  There is at most
one final inert switch.  Hence in the quadratic case the total
number $t$ of conductors satisfies, for example,
\begin{equation}\label{eq:uniform-conductor-count}
 t\le 2+C_B+4b+(1+4b)(C_{\QQ}+4).
\end{equation}
For $B=\QQ$ the direct loop gives the simpler bound
$t\le1+C_{\QQ}+4$.  These constants depend only on the fixed
original data and the fixed choice of $B$.

Every original $\ell\in\Sigma$ has full Frobenius order in $L$ and
its degree-$2^a$ layer, so has at most $b$ primes in either $A$ or
$F$.  Every conductor is totally ramified in both cyclic layers,
so likewise has at most $b$ primes.  Consequently
\[
                         u_A,u_F\le b(|\Sigma|+t).
\]
The analogous scalar bound is $2|\Sigma|$, and the absolutely
irreducible bound is $|\Sigma|$.  This proves the required common
constant $U$.  The degrees are $d_A=b2^a$, $d=b2^n$, and
$m=2^{n-a}$ in the marked cases.  Taking $a$ and $n-a$ large makes
$d_A$ and $[F:A]$ large independently.  In the marked cases this
makes $m$ large and gives $m/d=1/d_A$.

Since $K_\alpha/\QQ$ has odd degree, $\alpha$ remains nontrivial
on every field constructed in the distinct-character case.
Inflation--restriction, with
$H^0(G_F,k(\beta))=0$, proves injectivity of restriction from $A$
to $F$.  Together with the scalar and absolutely irreducible cases
proved above, this establishes all the field and restriction
assertions of Proposition~\ref{prop:field-preparation}.

It remains to derive \eqref{eq:strict-selmer}.  Since the mod-two
cyclotomic character is trivial, the Tate dual of $k(\beta)$ is
$k(\beta^{-1})$.  Its Selmer conditions dual to the strict ones
are exactly those in \eqref{eq:dual-vanishing}.  The dual Selmer
group is zero.  The global invariant terms in the Selmer dimension
formula cancel, whether $\beta=1$ or $\beta\ne1$.  At a real place,
the full condition contributes
$\dim H^1_v-h^0_v=1-1=0$.  At an unmarked dyadic place, local
Euler characteristic gives contribution one if $\alpha_2\ne1$
and two if $\alpha_2=1$.  At a marked place the zero condition
contributes $-h^0_v=-1$.  At a non-dyadic allowed place the full
condition contributes $h^2_v\le1$, and an unramified condition
elsewhere contributes zero.  Adding these terms gives exactly the
two bounds in \eqref{eq:strict-selmer}.

For later use, the degree choices can in particular ensure the strict
inequalities
\begin{align}
 d+u_F+3&<2d-1 &&\text{in the unmarked reducible case},\label{eq:field-unmarked-gap}\\
 2d-2m+u_F+3&<2d-m-1 &&\text{in the marked case},\label{eq:field-marked-gap}\\
 2d-2m-4-u_F&>d.\label{eq:field-curve-gap}
\end{align}
For example, it is enough in the marked case to have
$m>U+4$ and $m(d_A-2)>U+4$.  The same freedom makes $d_A$
larger than any prescribed fixed cost in the initial trace cuts.
These choices are legitimate because \eqref{eq:uniform-conductor-count}
and the bounds for $u_A,u_F$ are independent of both cyclic exponents.
